% GENERATED FILE. Edit canonical lessons, then run npm run generate:books.
% canonical-source-hash: fnv1a64:7792d9f0818df951
% canonical-lessons: ES-C612-a-no-ser-que, ES-C612-a-condicion-de-que, ES-C612-de-lo-contrario, ES-C612-en-ese-caso, ES-C612-en-todo-caso

\chapter{Si no es así}
\label{ch:es-si-no-es-asi}
\hlchaptermodality{612}

\begin{chapteropening}
\textbf{By the end of this chapter:} \emph{I can say what happens otherwise, in that case or in any case.}

Complete the last lesson of chapter 612: I can say what happens otherwise, in that case or in any case.
\end{chapteropening}

\section[a no ser que]{a no ser que — except if; unless}

\label{lesson:ES-C612-a-no-ser-que}



\begin{quote}
Before the new one: say the Spanish for provided that, then the Spanish for in case.
\end{quote}

\subsection*{You'll want to know: a no ser que}
\textbf{a no ser que} — \textquotedblleft{}except if; unless\textquotedblright{}.

\begin{quote}
Salimos a las ocho, a no ser que llueva.
\end{quote}

\begin{quote}
\emph{We leave at eight, unless it rains.}
\end{quote}

Means the same as \emph{a menos que}, with the subjunctive.

\begin{cousinweb}
\emph{a no ser que} is built on \emph{ser}, usually explained as a merger of Latin \emph{esse} 'to be' and \emph{sedere} 'to sit'; EtymDB takes it back to Old Spanish \emph{seer}. The phrase means 'unless it be that'.
\end{cousinweb}

\begin{grammarlens}[title={What you've built}]
One more word for telling, explaining and describing at B1.
\end{grammarlens}

\subsection*{Your turn}
\begin{itemize}
\raggedright
  \item \emph{Say it:} \emph{a no ser que}
  \item \emph{Say it:} \emph{a no ser que}, once more
  \item \emph{Say it:} a sentence of your own with \emph{a no ser que}
  \item \emph{Recall:} say the Spanish for provided that, then the Spanish for in case, then say \emph{a no ser que} again
\end{itemize}

\begin{tcolorbox}[breakable,colback=teal!4,colframe=teal!35!black,title={Before you move on}]
What does a no ser que mean? (\textquotedblleft{}Except if; unless\textquotedblright{}.) Say it once more.
\end{tcolorbox}
\section[a condición de que]{a condición de que — on condition that}

\label{lesson:ES-C612-a-condicion-de-que}



\begin{quote}
Before the new one: say the Spanish for in case, then the Spanish for except if.
\end{quote}

\subsection*{You'll want to know: a condición de que}
\textbf{a condición de que} — \textquotedblleft{}on condition that\textquotedblright{}.

\begin{quote}
Firmo a condición de que me paguen antes.
\end{quote}

\begin{quote}
\emph{I will sign on condition that they pay me first.}
\end{quote}

Formal; followed by the subjunctive.

\begin{cousinweb}
\emph{a condición de que} is built on \emph{condición}, a learned borrowing from Latin \emph{condicio} 'terms, agreement' (EtymDB).
\end{cousinweb}

\begin{grammarlens}[title={What you've built}]
One more word for telling, explaining and describing at B1.
\end{grammarlens}

\subsection*{Your turn}
\begin{itemize}
\raggedright
  \item \emph{Say it:} \emph{a condición de que}
  \item \emph{Say it:} \emph{a condición de que}, once more
  \item \emph{Say it:} a sentence of your own with \emph{a condición de que}
  \item \emph{Recall:} say the Spanish for in case, then the Spanish for except if, then say \emph{a condición de que} again
\end{itemize}

\begin{tcolorbox}[breakable,colback=teal!4,colframe=teal!35!black,title={Before you move on}]
What does a condición de que mean? (\textquotedblleft{}On condition that\textquotedblright{}.) Say it once more.
\end{tcolorbox}
\section[de lo contrario]{de lo contrario — otherwise}

\label{lesson:ES-C612-de-lo-contrario}



\begin{quote}
Before the new one: say the Spanish for except if, then the Spanish for on condition that.
\end{quote}

\subsection*{You'll want to know: de lo contrario}
\textbf{de lo contrario} — \textquotedblleft{}otherwise\textquotedblright{}.

\begin{quote}
Date prisa; de lo contrario, perdemos el tren.
\end{quote}

\begin{quote}
\emph{Hurry up; otherwise, we will miss the train.}
\end{quote}

In speech people also say \emph{si no}.

\begin{cousinweb}
\emph{de lo contrario} is built on \emph{contrario}, a learned borrowing from Latin \emph{contrarius} 'opposite', from \emph{contra} 'against' (Lewis \& Short).
\end{cousinweb}

\begin{grammarlens}[title={What you've built}]
One more word for telling, explaining and describing at B1.
\end{grammarlens}

\subsection*{Your turn}
\begin{itemize}
\raggedright
  \item \emph{Say it:} \emph{de lo contrario}
  \item \emph{Say it:} \emph{de lo contrario}, once more
  \item \emph{Say it:} a sentence of your own with \emph{de lo contrario}
  \item \emph{Recall:} say the Spanish for except if, then the Spanish for on condition that, then say \emph{de lo contrario} again
\end{itemize}

\begin{tcolorbox}[breakable,colback=teal!4,colframe=teal!35!black,title={Before you move on}]
What does de lo contrario mean? (\textquotedblleft{}Otherwise\textquotedblright{}.) Say it once more.
\end{tcolorbox}
\section[en ese caso]{en ese caso — in that case}

\label{lesson:ES-C612-en-ese-caso}



\begin{quote}
Before the new one: say the Spanish for on condition that, then the Spanish for otherwise.
\end{quote}

\subsection*{You'll want to know: en ese caso}
\textbf{en ese caso} — \textquotedblleft{}in that case\textquotedblright{}.

\begin{quote}
—No tengo coche. —En ese caso, vamos en tren.
\end{quote}

\begin{quote}
\emph{I do not have a car. In that case, we will go by train.}
\end{quote}

It answers a condition that someone has given.

\begin{cousinweb}
\emph{en ese caso} is built on \emph{caso}, a learned borrowing from Latin \emph{casus} 'an event', from \emph{cadere} 'to fall' (EtymDB).
\end{cousinweb}

\begin{grammarlens}[title={What you've built}]
One more word for telling, explaining and describing at B1.
\end{grammarlens}

\subsection*{Your turn}
\begin{itemize}
\raggedright
  \item \emph{Say it:} \emph{en ese caso}
  \item \emph{Say it:} \emph{en ese caso}, once more
  \item \emph{Say it:} a sentence of your own with \emph{en ese caso}
  \item \emph{Recall:} say the Spanish for on condition that, then the Spanish for otherwise, then say \emph{en ese caso} again
\end{itemize}

\begin{tcolorbox}[breakable,colback=teal!4,colframe=teal!35!black,title={Before you move on}]
What does en ese caso mean? (\textquotedblleft{}In that case\textquotedblright{}.) Say it once more.
\end{tcolorbox}
\section[en todo caso]{en todo caso — in any case; if anything}

\label{lesson:ES-C612-en-todo-caso}



\begin{quote}
Before the new one: say the Spanish for otherwise, then the Spanish for in that case.
\end{quote}

\subsection*{You'll want to know: en todo caso}
\textbf{en todo caso} — \textquotedblleft{}in any case; if anything\textquotedblright{}.

\begin{quote}
En todo caso, te llamo mañana.
\end{quote}

\begin{quote}
\emph{In any case, I will call you tomorrow.}
\end{quote}

Close to \emph{de todas formas}.

\begin{cousinweb}
\emph{en todo caso} is built on \emph{todo}, inherited from Latin \emph{totus} 'whole', and \emph{caso}, borrowed from Latin \emph{casus} 'an event' (EtymDB).
\end{cousinweb}

\begin{grammarlens}[title={What you've built}]
One more word for telling, explaining and describing at B1.
\end{grammarlens}

\subsection*{Your turn}
\begin{itemize}
\raggedright
  \item \emph{Say it:} \emph{en todo caso}
  \item \emph{Say it:} \emph{en todo caso}, once more
  \item \emph{Say it:} a sentence of your own with \emph{en todo caso}
  \item \emph{Recall:} say the Spanish for otherwise, then the Spanish for in that case, then say \emph{en todo caso} again
\end{itemize}

\begin{tcolorbox}[breakable,colback=teal!4,colframe=teal!35!black,title={Before you move on}]
What does en todo caso mean? (\textquotedblleft{}In any case; if anything\textquotedblright{}.) Say it once more.
\end{tcolorbox}
